% generated by docs/audit/make_tables.py -- do not edit by hand
\begingroup
\setlength{\tabcolsep}{4pt}
\renewcommand{\arraystretch}{1.15}
\noindent\begin{tabular}{@{}p{0.27\textwidth}@{\hspace{9pt}}p{0.70\textwidth}@{}}
\toprule
\textbf{the mathematics} & \textbf{the Lean, verbatim}\\
\bottomrule
\end{tabular}\par\smallskip
\noindent\begin{tabular}{@{}p{0.27\textwidth}@{\hspace{9pt}}p{0.70\textwidth}@{}}
\textbf{(1)}\ $P$ is a closed heptagon that does not cross itself: no side is a point, non-adjacent sides are disjoint, and two consecutive sides meet exactly in the vertex they share. & {\scriptsize\ttfamily\setlength{\parindent}{0pt}\raggedright
\mbox{}def~SimpleClosed~(W~:~Fin~7~$\to$~$\mathbb{R}$~$\times$~$\mathbb{R}$)~:~Prop~:=\\[-1pt]
\mbox{}~~($\forall$~i~:~Fin~7,~W~i~$\ne$~W~(i~+~1))~$\land$\\[-1pt]
\mbox{}~~($\forall$~i~j~:~Fin~7,~$\lnot$~LowerBound.Adjacent~i~j~$\to$\\[-1pt]
\mbox{}~~~~~~Disjoint~(segment~$\mathbb{R}$~(W~i)~(W~(i~+~1)))~(segment~$\mathbb{R}$~(W~j)~(W~(j~+~1))))~$\land$\\[-1pt]
\mbox{}~~($\forall$~i~:~Fin~7,\\[-1pt]
\mbox{}~~~~~~segment~$\mathbb{R}$~(W~i)~(W~(i~+~1))~$\cap$~segment~$\mathbb{R}$~(W~(i~+~1))~(W~(i~+~2))\\[-1pt]
\mbox{}~~~~~~~~=~\{W~(i~+~1)\})\\[-1pt]
}\par\smallskip{\footnotesize\ttfamily lean/AKR9.lean:71}\\
\end{tabular}
\par\smallskip\noindent\hspace*{0.29\textwidth}\begin{minipage}{0.70\textwidth}{\footnotesize\textit{Note.}\ This is simplicity \emph{as a curve} and nothing more. Three collinear vertices are deliberately allowed --- a straight vertex is not excluded here. \texttt{Adjacent} is cyclic adjacency on $\mathbb{Z}/7$.}\end{minipage}
\par\medskip\noindent\rule{\textwidth}{0.3pt}\par\smallskip
\noindent\begin{tabular}{@{}p{0.27\textwidth}@{\hspace{9pt}}p{0.70\textwidth}@{}}
\textbf{(2)}\ The object under study is an ordered pair $(P,Q)$ of such heptagons in the plane. Nothing else is assumed about it. & {\scriptsize\ttfamily\setlength{\parindent}{0pt}\raggedright
\mbox{}structure~Pair9~where\\[-1pt]
\mbox{}~~/--~The~fourteen~vertices:~`P`~is~`0..6`,~`Q`~is~`7..13`.~-/\\[-1pt]
\mbox{}~~V~:~Fin~14~$\to$~$\mathbb{R}$~$\times$~$\mathbb{R}$\\[-1pt]
\mbox{}~~simpleP~:~SimpleClosed~(fun~i~:~Fin~7~=>~V~(pP~i))\\[-1pt]
\mbox{}~~simpleQ~:~SimpleClosed~(fun~i~:~Fin~7~=>~V~(pQ~i))\\[-1pt]
}\par\smallskip{\footnotesize\ttfamily lean/AKR9.lean:81}\\
\end{tabular}
\par\smallskip\noindent\hspace*{0.29\textwidth}\begin{minipage}{0.70\textwidth}{\footnotesize\textit{Note.}\ The fourteen vertices are one function \texttt{V}; \texttt{pP\,i} $=i$ and \texttt{pQ\,i} $=7+i$ split them. There is no general-position field, and no non-degeneracy field. The two \texttt{simple} fields are the whole hypothesis.}\end{minipage}
\par\medskip\noindent\rule{\textwidth}{0.3pt}\par\smallskip
\noindent\begin{tabular}{@{}p{0.27\textwidth}@{\hspace{9pt}}p{0.70\textwidth}@{}}
\textbf{(3)}\ $\partial P=\bigcup_{i\in\mathbb{Z}/7}[P_i,P_{i+1}]$, the union of the seven closed sides; likewise $\partial Q$. & {\scriptsize\ttfamily\setlength{\parindent}{0pt}\raggedright
\mbox{}def~bdryP9~(R~:~Pair9)~:~Set~($\mathbb{R}$~$\times$~$\mathbb{R}$)~:=~$\bigcup$~i~:~Fin~7,~sideP9~R~i\\[-1pt]
}\par\smallskip{\footnotesize\ttfamily lean/AKR9.lean:115}\\
\end{tabular}
\par\smallskip\noindent\hspace*{0.29\textwidth}\begin{minipage}{0.70\textwidth}{\footnotesize\textit{Note.}\ \texttt{segment $\mathbb{R}$ a b} is Mathlib's \emph{closed} segment $\{(1-t)a+tb : 0\le t\le 1\}$, so endpoints are included. \texttt{sideP9} is that segment; \texttt{bdryQ9} is the same with \texttt{pQ}.}\end{minipage}
\par\medskip\noindent\rule{\textwidth}{0.3pt}\par\smallskip
\noindent\begin{tabular}{@{}p{0.27\textwidth}@{\hspace{9pt}}p{0.70\textwidth}@{}}
\textbf{(4)}\ The set being counted is $\partial P\cap\partial Q$: the points of the plane lying on both boundaries. & {\scriptsize\ttfamily\setlength{\parindent}{0pt}\raggedright
\mbox{}def~points9~(R~:~Pair9)~:~Set~($\mathbb{R}$~$\times$~$\mathbb{R}$)~:=~bdryP9~R~$\cap$~bdryQ9~R\\[-1pt]
}\par\smallskip{\footnotesize\ttfamily lean/AKR9.lean:119}\\
\end{tabular}
\par\smallskip\noindent\hspace*{0.29\textwidth}\begin{minipage}{0.70\textwidth}{\footnotesize\textit{Note.}\ A set of points, not a list of crossings and not a multiset. A point where three sides concur is one element, not two.}\end{minipage}
\par\medskip\noindent\rule{\textwidth}{0.3pt}\par\smallskip
\noindent\begin{tabular}{@{}p{0.27\textwidth}@{\hspace{9pt}}p{0.70\textwidth}@{}}
\textbf{(5)}\ \emph{Assumed:} no vertex of $P$ equals a vertex of $Q$. (This is the conclusion of AKR \S9.1; see \S7.) & {\scriptsize\ttfamily\setlength{\parindent}{0pt}\raggedright
\mbox{}def~attainedCoincFree~:~Set~$\mathbb{N}$~:=\\[-1pt]
\mbox{}~~\{~n~|~$\exists$~R~:~Pair9,~(points9~R).Finite~$\land$\\[-1pt]
\mbox{}~~~~~~($\forall$~v~j~:~Fin~7,~R.V~(pP~v)~$\ne$~R.V~(pQ~j))~$\land$~(points9~R).ncard~=~n~\}\\[-1pt]
}\par\smallskip{\footnotesize\ttfamily lean/Pair9Shift.lean:4314}\\
\end{tabular}
\par\smallskip\noindent\hspace*{0.29\textwidth}\begin{minipage}{0.70\textwidth}{\footnotesize\textit{Note.}\ The middle clause is the assumption, written out. It forbids \emph{equality} of vertices only --- a vertex lying in the interior of the other polygon's side is \emph{not} excluded, and is handled by the proof.}\end{minipage}
\par\medskip\noindent\rule{\textwidth}{0.3pt}\par\smallskip
\noindent\begin{tabular}{@{}p{0.27\textwidth}@{\hspace{9pt}}p{0.70\textwidth}@{}}
\textbf{(6)}\ If $\partial P\cap\partial Q$ is finite and $P,Q$ share no vertex, then $\#(\partial P\cap\partial Q)\le 38$. & {\scriptsize\ttfamily\setlength{\parindent}{0pt}\raggedright
\mbox{}theorem~points9\_le\_38\_of\_coincFree~(R~:~Pair9)~(hfin~:~(points9~R).Finite)\\[-1pt]
\mbox{}~~~~(hcoinc~:~$\forall$~v~j~:~Fin~7,~R.V~(pP~v)~$\ne$~R.V~(pQ~j))~:\\[-1pt]
\mbox{}~~~~(points9~R).ncard~$\le$~38~:=~by\\[-1pt]
\mbox{}~~obtain~$\langle$Rend,~hD,~hE,~hb$\rangle$~:=~exists\_chain\_deg\_free~\_~R~hfin~hcoinc~rfl\\[-1pt]
\mbox{}~~exact~points9\_le\_38\_of\_joint~R~Rend~hfin~hb~hD~hE\\[-1pt]
}\par\smallskip{\footnotesize\ttfamily lean/Pair9Shift.lean:4223}\\
\end{tabular}
\par\smallskip\noindent\hspace*{0.29\textwidth}\begin{minipage}{0.70\textwidth}{\footnotesize\textit{Note.}\ \texttt{ncard} is the cardinality of a set as a natural number. \texttt{hfin} is the finiteness hypothesis; without it the count is not a number at all (see the last row of \S3).}\end{minipage}
\par\medskip\noindent\rule{\textwidth}{0.3pt}\par\smallskip
\noindent\begin{tabular}{@{}p{0.27\textwidth}@{\hspace{9pt}}p{0.70\textwidth}@{}}
\textbf{(7)}\ $38$ is attained, and is the largest attained value: it is the maximum, not merely an upper bound. & {\scriptsize\ttfamily\setlength{\parindent}{0pt}\raggedright
\mbox{}theorem~isGreatest\_38\_of\_coincFree~:~IsGreatest~attainedCoincFree~38~:=~by\\[-1pt]
\mbox{}~~refine~$\langle$$\langle$AKR9.certPair9,~certPair9\_finite,~certPair9\_coincFree,\\[-1pt]
\mbox{}~~~~AKR9.certPair9\_ncard$\rangle$,~?\_$\rangle$\\[-1pt]
\mbox{}~~rintro~n~$\langle$R,~hfin,~hco,~rfl$\rangle$\\[-1pt]
\mbox{}~~exact~points9\_le\_38\_of\_coincFree~R~hfin~hco\\[-1pt]
}\par\smallskip{\footnotesize\ttfamily lean/Pair9Shift.lean:4324}\\
\end{tabular}
\par\smallskip\noindent\hspace*{0.29\textwidth}\begin{minipage}{0.70\textwidth}{\footnotesize\textit{Note.}\ \texttt{IsGreatest S a} unfolds to $a\in S \wedge \forall b\in S,\ b\le a$. The membership half is the explicit certificate of \S3; the bound half is the row above.}\end{minipage}
\par\medskip\noindent\rule{\textwidth}{0.3pt}\par\smallskip
\endgroup
